% Generated by scripts/assemble.py; corrections belong in source/corrections.json.
\examhead{177}{西北大学 2025 年数学分析试卷}{167}
% p168-b01; source PDF page 168
\blockanchor{p168-b01}
\question{1}求极限.\par
% p168-b02; source PDF page 168
\blockanchor{p168-b02}
(1) $\lim_{x\to 0}\frac{\cos\left(\frac{\pi}{2}\cos x\right)}{\sin(\sin x)}.$\par
% p168-b03; source PDF page 168
\blockanchor{p168-b03}
(2) $\lim_{n\to \infty}\int_0^1 x^n\sqrt{1 + x^2}\mathrm{d}x.$\par
% p168-b04; source PDF page 168
\blockanchor{p168-b04}
\question{2}已知数列 $\{x_{n}\}$ 满足 $x_{1} = 2, x_{n+1} = \frac{1}{2}\left(x_{n} + \frac{1}{x_{n}}\right), n = 1, 2, \cdots$.\par
% p168-b05; source PDF page 168
\blockanchor{p168-b05}
(1) 证明: $\lim_{n\to \infty}x_n$ 存在并求其值.\par
% p168-b06; source PDF page 168
\blockanchor{p168-b06}
(2) 证明: 级数 $\sum_{n=1}^{\infty}\left(\frac{x_n}{x_{n+1}} - 1\right)$ 收敛.\par
% p168-b07; source PDF page 168
\par\Needspace{6\baselineskip}
\blockanchor{p168-b07}
\question{3}已知 $f(x)$ 在 $[0,1]$ 上可导，且 $f(1) = \frac{1}{\alpha} \int_{0}^{\alpha} xe^{1 - x} f(x) \, \mathrm{d}x$，证明：存在 $\xi \in (0,1)$，使得\par
\nopagebreak[4]
% p168-b08; source PDF page 168
\blockanchor{p168-b08}
\displayformula{f ^ {\prime} (\xi) = \left(1 - \frac {1}{\xi}\right) f (\xi).}
% p168-b09; source PDF page 168
\blockanchor{p168-b09}
\question{4}已知 $a > 0$, 证明: 曲面 $\sqrt{x} + \sqrt{y} + \sqrt{z} = \sqrt{a}$ 上任一点的切平面在各坐标轴上的截距之和均为 $a$.\par
% p168-b10; source PDF page 168
\blockanchor{p168-b10}
\question{5}已知 $n$ 为正整数, $f(x) = \sum_{n=1}^{\infty} \frac{\sin nx}{n+1}$.\par
% p168-b11; source PDF page 168
\blockanchor{p168-b11}
(1) 证明: $f(x)$ 在 $(0,2\pi)$ 上连续.\par
% p168-b12; source PDF page 168
\blockanchor{p168-b12}
(2) 求幂级数 $\sum_{n=1}^{\infty} (-1)^{n-1} \frac{x^n}{n}$ 的和函数.\par
% p168-b13; source PDF page 168
\blockanchor{p168-b13}
(3) 证明: $\int_{0}^{\pi} f(x) \, \mathrm{d}x = 2 \ln 2$.\par
% p168-b14; source PDF page 168
\par\Needspace{6\baselineskip}
\blockanchor{p168-b14}
\question{6}计算三重积分\par
\nopagebreak[4]
% p168-b15; source PDF page 168
\blockanchor{p168-b15}
\displayformula{\iiint_ {\Omega} (x + z) \mathrm{d} x \mathrm{d} y \mathrm{d} z.}
% p168-b16; source PDF page 168
\blockanchor{p168-b16}
其中 $\Omega$ 为锥面 $z = \sqrt{x^2 + y^2}$ 与上半球面 $z = \sqrt{1 - x^2 - y^2}$ 围成的区域.\par
% p168-b17; source PDF page 168
\par\Needspace{6\baselineskip}
\blockanchor{p168-b17}
\question{7}求第一型曲线积分\par
\nopagebreak[4]
% p168-b18; source PDF page 168
\blockanchor{p168-b18}
\displayformula{\int_ {L} \left(x ^ {2} + 2 y ^ {2} + 3 z\right) \mathrm{d} s.}
% p168-b19; source PDF page 168
\blockanchor{p168-b19}
其中 $L$ 是 $x^{2} + y^{2} + z^{2} = R^{2}(R > 0)$ 和 $x + y + z = 0$ 的交线.\par
% p168-b20; source PDF page 168
\par\Needspace{6\baselineskip}
\blockanchor{p168-b20}
\question{8}计算第二类曲面积分\par
\nopagebreak[4]
% p168-b21; source PDF page 168
\blockanchor{p168-b21}
\displayformula{\iint_ {\Sigma} \frac {x \mathrm{d} y \mathrm{d} z - y \mathrm{d} z \mathrm{d} x + (z ^ {3} + 1) \mathrm{d} x \mathrm{d} y}{(2 x ^ {2} + 3 y ^ {2} + z ^ {2}) ^ {3}}.}
% p168-b22; source PDF page 168
\blockanchor{p168-b22}
其中 $\Sigma$ 为 $z = \sqrt{1 - 2x^2 - 3y^2}$ 的上侧.\par
% p168-b23; source PDF page 168
\par\Needspace{6\baselineskip}
\blockanchor{p168-b23}
\question{9}设 $\varphi(t)$ 在 $[0,1]$ 上连续, 令\par
\nopagebreak[4]
% p168-b24; source PDF page 168
\blockanchor{p168-b24}
\displayformula{f (x) = \int_ {0} ^ {1} \frac {x}{x ^ {2} + t ^ {2}} \varphi (t) \mathrm{d} t.}
% p168-b25; source PDF page 168
\blockanchor{p168-b25}
证明： $f(x)$ 在 $x = 0$ 处连续的充分必要条件是 $\varphi (0) = 0$\par
